\documentclass[12pt]{article}

% Idioma y codificación
\usepackage[spanish, es-tabla]{babel}       %es-tabla para que se titule "Tabla"
\usepackage[utf8]{inputenc}

% Márgenes
\usepackage[a4paper,top=3cm,bottom=2.5cm,left=3cm,right=3cm]{geometry}

% Comentarios de bloque
\usepackage{verbatim}

% Paquetes de links
\usepackage[hidelinks]{hyperref}    % Permite enlaces
\usepackage{url}                    % redirecciona a la web

% Más opciones para enumeraciones
\usepackage{enumitem}

% Personalizar la portada
\usepackage{titling}


% Paquetes de tablas
\usepackage{multirow}


%------------------------------------------------------------------------

%Paquetes de figuras
\usepackage{caption}
\usepackage{subcaption} % Figuras al lado de otras
\usepackage{float}      % Poner figuras en el sitio indicado H.


% Paquetes de imágenes
\usepackage{graphicx}       % Paquete para añadir imágenes
\usepackage{transparent}    % Para manejar la opacidad de las figuras

% Paquete para usar colores
\usepackage[dvipsnames]{xcolor}
\usepackage{pagecolor}      % Para cambiar el color de la página

% Habilita tamaños de fuente mayores
\usepackage{fix-cm}


%------------------------------------------------------------------------

% Paquetes de matemáticas
\usepackage{mathtools, amsfonts, amssymb, mathrsfs}
\usepackage[makeroom]{cancel}     % Simplificar tachando
\usepackage{polynom}    % Divisiones y Ruffini
\usepackage{units} % Para poner fracciones diagonales con \nicefrac

\usepackage{pgfplots}   %Representar funciones
\pgfplotsset{compat=1.18}  % Versión 1.18

\usepackage{tikz-cd}    % Para usar diagramas de composiciones
\usetikzlibrary{calc}   % Para usar cálculo de coordenadas en tikz

% Paquete para usar el símbolo de grados
\usepackage{gensymb}

%Definición de teoremas, etc.
\usepackage{amsthm}
%\swapnumbers   % Intercambia la posición del texto y de la numeración

\theoremstyle{plain}

\makeatletter
\@ifclassloaded{article}{
  \newtheorem{teo}{Teorema}[section]
}{
  \theoremstyle{definition}
  \newtheorem{teo}{Teorema}[chapter]  % Se resetea en cada chapter
}
\makeatother

\theoremstyle{definition}
\newtheorem{coro}{Corolario}[teo]           % Se resetea en cada teorema
\newtheorem{prop}[teo]{Proposición}         % Usa el mismo contador que teorema
\newtheorem{lema}[teo]{Lema}                % Usa el mismo contador que teorema

\theoremstyle{remark}
\newtheorem*{observacion}{Observación}

\theoremstyle{definition}

\makeatletter
\@ifclassloaded{article}{
  \newtheorem{definicion}{Definición} [section]     % Se resetea en cada chapter
}{
  \newtheorem{definicion}{Definición} [chapter]     % Se resetea en cada chapter
}
\makeatother

\newtheorem*{notacion}{Notación}
\newtheorem*{ejemplo}{Ejemplo}
\newtheorem*{propiedades}{Propiedades}
\newtheorem*{ejercicio*}{Ejercicio}             % No numerado
\newtheorem{ejercicio}{Ejercicio} [section]     % Se resetea en cada section


% Modificar el formato de la numeración del teorema "ejercicio"
\renewcommand{\theejercicio}{%
  \ifnum\value{section}=0 % Si no se ha iniciado ninguna sección
    \arabic{ejercicio}% Solo mostrar el número de ejercicio
  \else
    \thesection.\arabic{ejercicio}% Mostrar número de sección y número de ejercicio
  \fi
}

% Símbolo de restricción
\newcommand\restr[2]{{% we make the whole thing an ordinary symbol
  \left.{#1}\right|_{#2}%
}}


\renewcommand\qedsymbol{\small{$\square$}}         % Cambiar símbolo QED
%------------------------------------------------------------------------

% Paquetes para encabezados
\usepackage{fancyhdr}
\pagestyle{fancy}
\fancyhf{}

\newcommand{\helv}{ % Modificación tamaño de letra
\fontfamily{}\fontsize{12}{12}\selectfont}
\setlength{\headheight}{15pt} % Amplía el tamaño del índice


%\usepackage{lastpage}   % Referenciar última pag   \pageref{LastPage}
\fancyfoot[C]{\thepage}

%------------------------------------------------------------------------

% Conseguir que no ponga "Capítulo 1". Sino solo "1."
\makeatletter
\@ifclassloaded{book}{
  \renewcommand{\chaptermark}[1]{\markboth{\thechapter.\ #1}{}} % En el encabezado
    
  \renewcommand{\@makechapterhead}[1]{%
  \vspace*{50\p@}%
  {\parindent \z@ \raggedright \normalfont
    \ifnum \c@secnumdepth >\m@ne
      \huge\bfseries \thechapter.\hspace{1em}\ignorespaces
    \fi
    \interlinepenalty\@M
    \Huge \bfseries #1\par\nobreak
    \vskip 40\p@
  }}
}
\makeatother

%------------------------------------------------------------------------
% Paquetes de cógido
\usepackage{minted}
\renewcommand\listingscaption{Código fuente}

\usepackage{fancyvrb}
% Personaliza el tamaño de los números de línea
\renewcommand{\theFancyVerbLine}{\small\arabic{FancyVerbLine}}

% Estilo para C++
\newminted{cpp}{
    frame=lines,
    framesep=2mm,
    baselinestretch=1.2,
    linenos,
    escapeinside=||
}


\usepackage{listings} % Para incluir código desde un archivo

\renewcommand\lstlistingname{Código Fuente}
\renewcommand\lstlistlistingname{Índice de Códigos Fuente}

% Definir colores
\definecolor{vscodepurple}{rgb}{0.5,0,0.5}
\definecolor{vscodeblue}{rgb}{0,0,0.8}
\definecolor{vscodegreen}{rgb}{0,0.5,0}
\definecolor{vscodegray}{rgb}{0.5,0.5,0.5}
\definecolor{vscodebackground}{rgb}{0.97,0.97,0.97}
\definecolor{vscodelightgray}{rgb}{0.9,0.9,0.9}

% Configuración para el estilo de C similar a VSCode
\lstdefinestyle{vscode_C}{
  backgroundcolor=\color{vscodebackground},
  commentstyle=\color{vscodegreen},
  keywordstyle=\color{vscodeblue},
  numberstyle=\tiny\color{vscodegray},
  stringstyle=\color{vscodepurple},
  basicstyle=\scriptsize\ttfamily,
  breakatwhitespace=false,
  breaklines=true,
  captionpos=b,
  keepspaces=true,
  numbers=left,
  numbersep=5pt,
  showspaces=false,
  showstringspaces=false,
  showtabs=false,
  tabsize=2,
  frame=tb,
  framerule=0pt,
  aboveskip=10pt,
  belowskip=10pt,
  xleftmargin=10pt,
  xrightmargin=10pt,
  framexleftmargin=10pt,
  framexrightmargin=10pt,
  framesep=0pt,
  rulecolor=\color{vscodelightgray},
  backgroundcolor=\color{vscodebackground},
}

%------------------------------------------------------------------------

% Comandos definidos
\newcommand{\bb}[1]{\mathbb{#1}}
\newcommand{\cc}[1]{\mathcal{#1}}

% I prefer the slanted \leq
\let\oldleq\leq % save them in case they're every wanted
\let\oldgeq\geq
\renewcommand{\leq}{\leqslant}
\renewcommand{\geq}{\geqslant}

% Si y solo si
\newcommand{\sii}{\iff}

% Letras griegas
\newcommand{\eps}{\epsilon}
\newcommand{\veps}{\varepsilon}
\newcommand{\lm}{\lambda}

\newcommand{\ol}{\overline}
\newcommand{\ul}{\underline}
\newcommand{\wt}{\widetilde}
\newcommand{\wh}{\widehat}

\let\oldvec\vec
\renewcommand{\vec}{\overrightarrow}

% Derivadas parciales
\newcommand{\del}[2]{\frac{\partial #1}{\partial #2}}
\newcommand{\Del}[3]{\frac{\partial^{#1} #2}{\partial^{#1} #3}}
\newcommand{\deld}[2]{\dfrac{\partial #1}{\partial #2}}
\newcommand{\Deld}[3]{\dfrac{\partial^{#1} #2}{\partial^{#1} #3}}


\newcommand{\AstIg}{\stackrel{(\ast)}{=}}
\newcommand{\Hop}{\stackrel{L'H\hat{o}pital}{=}}

\newcommand{\red}[1]{{\color{red}#1}} % Para integrales, destacar los cambios.

% Método de integración
\newcommand{\MetInt}[2]{
    \left[\begin{array}{c}
        #1 \\ #2
    \end{array}\right]
}

% Declarar aplicaciones
% 1. Nombre aplicación
% 2. Dominio
% 3. Codominio
% 4. Variable
% 5. Imagen de la variable
\newcommand{\Func}[5]{
    \begin{equation*}
        \begin{array}{rrll}
            #1:& #2 & \longrightarrow & #3\\
               & #4 & \longmapsto & #5
        \end{array}
    \end{equation*}
}

%------------------------------------------------------------------------
\setlength{\parindent}{0pt} % quitar la sangría

\setlist[itemize, 1]{label=\textbullet \textbf{)}}  % Cambiar las viñetas a círculos

\def\endsquare{%
  \begin{flushright}%
      \small{$\square$}%
  \end{flushright}%
}

\newcommand{\apuntar}[1]{%
  \hyperlink{#1}{\textbf{(#1)}}%p
}

\newcommand{\objetivo}[1]{%
  \textbf{(\hypertarget{#1}{#1})}
}

\usepackage{stackrel}

\newcommand{\captext}[1]{\stackbin [#1]{}{\cap}}
\newcommand{\cuptext}[1]{\stackbin [#1]{}{\cup}}


\usepackage{xkeyval} % Para el paso de argumentos

\usepackage{graphicx}

% Definir la carpeta de las imágenes
\graphicspath{{../_assets}{../../_assets}{../../../_assets}{../../../../_assets}{../../../../../_assets}}

% Definir el comando \portada
\makeatletter
\define@key{portada}{titulo}{\def\titulo{#1}}
\define@key{portada}{subtitulo}{\def\subtitulo{#1}}
\define@key{portada}{autor}{\def\autor{#1}}
\define@key{portada}{año}{\def\año{#1}}
\define@key{portada}{color página}{\def\colorpagina{#1}}
\define@key{portada}{imagen}{\def\imagen{#1}}

\newcommand*{\portada}[1][]{%
  % Encabezado
  \fancyhead[L]{\helv \titulo}
  \fancyhead[R]{\helv \nouppercase{\leftmark}}

    % Definimos las claves y sus valores por defecto
  \setkeys{portada}{%
    titulo=Sin Título,%
    subtitulo=Sin Subtítulo,%
    autor=Autor Desconocido,%
    año=Sin Año,%
    color página=white,%
    imagen=Logo-UGR-Black.png, #1}%
    \begin{titlepage}
        \centering
        \pagecolor{\colorpagina} % Cambia el color de fondo de la portada
        {\includegraphics[width=0.2\textwidth]{\imagen}\par}
        \vspace{1cm}
        {\bfseries\LARGE Universidad de Granada \par}
        \vspace{1cm}
        {\scshape\Large Doble Grado en Ingeniería Informática y Matemáticas \par}
        \vspace{3cm}
        {\scshape\Huge \titulo \par}
        \vspace{3cm}
        {\itshape\Large \subtitulo \par}
        \vfill
        {\Large Autor: \par}
        {\Large \autor \par}
        \vfill
        {\Large \año \par}
    \end{titlepage}%
    \pagecolor{white}
}


\usepackage[all]{xy}
\usepackage{ stmaryrd }
\usepackage{pgfplots}


\newcommand{\N}{\mathcal{N}}
\newcommand{\R}{\mathbb{R}}
\renewcommand{\P}{{P}}

\begin{document}
    \portada[%
        titulo=Inferencia Estadística,
        subtitulo=Entrega Ejercicio Tema 2,
        autor=Jesús Muñoz Velasco,
        año=Curso 2026-2027]
        
    \begin{ejercicio}
        Sean $(X_1,\dots,X_n)$ e $(Y_1,\dots,Y_m)$ muestras aleatorias simples independientes de poblaciones $\N(\mu_1,\sigma^2)$ y $\N(\mu_2,4\sigma^2)$, respectivamente. Sean $\alpha,\beta\in\R$ no ambos nulos, y sean $\overline{X}$, $\overline{Y}$, $S_X^2$, $S_Y^2$ las medias y cuasivarianzas de las dos muestras. Se pide:
    
        \begin{enumerate}[label=\alph*)]
            \item Calcular la distribución de
            \begin{gather*}
                T=\frac{\alpha(\overline{X}-\mu_1)+\beta(\overline{Y}-\mu_2)}{\sqrt{\dfrac{(n-1)S_X^2+\dfrac{(m-1)S_Y^2}{4}}{n+m-2}}\;\sqrt{\dfrac{\alpha^2}{n}+\dfrac{4\beta^2}{m}}}.
            \end{gather*}
            \item Para $n=4$ y $m=5$ calcular $\P[T \geq 2.365]$.
        \end{enumerate}
        \vspace{1cm}
        \subsection*{Resolución}

        \begin{enumerate}[label=\alph*)]
            \item Lo primero que observamos es que, por ser las muestras normales, por el Lema de Fisher tenemos que
            \begin{align*}
                &\overline{X}\sim\N\!\left(\mu_1,\frac{\sigma^2}{n}\right), &&\frac{(n-1)S_X^2}{\sigma^2}\sim\chi^2_{n-1}, &&\overline{X}\text{ y }S_X^2\text{ independientes},\\[2pt]
                &\overline{Y}\sim\N\!\left(\mu_2,\frac{4\sigma^2}{m}\right), &&\frac{(m-1)S_Y^2}{4\sigma^2}\sim\chi^2_{m-1}, &&\overline{Y}\text{ y }S_Y^2\text{ independientes}.
            \end{align*}
            Como además las dos muestras son independientes entre sí, las cuatro variables $\overline{X}$, $\overline{Y}$, $S_X^2$, $S_Y^2$ son independientes.\\

            Definimos
            \[
                U=\frac{(n-1)S_X^2+\dfrac{(m-1)S_Y^2}{4}}{\sigma^2}=\frac{(n-1)S_X^2}{\sigma^2}+\frac{(m-1)S_Y^2}{4\sigma^2}.
            \]
            El primer sumando es una $\chi^2_{n-1}$ y el segundo una $\chi^2_{m-1}$, y son independientes porque dependen de muestras independientes. Por la propiedad reproductiva de la chi-cuadrado,
            \[
                U\sim\chi^2_{n+m-2}.
            \]

            Denotamos ahora $W=\alpha(\overline{X}-\mu_1)+\beta(\overline{Y}-\mu_2)$. Es combinación lineal de normales independientes, luego es normal con
            \begin{align*}
                E[W]&=E[\alpha(\overline{X}-\mu_1)+\beta(\overline{Y}-\mu_2)] = \alpha(E[\overline{X}] - \mu_1) + \beta(E[\overline{Y}] - \mu_2)=\\
                &=\alpha(\mu_1-\mu_1)+\beta(\mu_2-\mu_2)=0\\
                \\
                Var(W)&= Var(\alpha(\overline{X}-\mu_1)+\beta(\overline{Y}-\mu_2)) = \alpha^2 Var(\overline{X}) + \beta^2 Var(\overline{Y})=\\
                &=\alpha^2\frac{\sigma^2}{n}+\beta^2\frac{4\sigma^2}{m}=\sigma^2\left(\frac{\alpha^2}{n}+\frac{4\beta^2}{m}\right).
            \end{align*}
            donde hemos usado que $\overline{X}$ e $\overline{Y}$ son independientes. Tipificando,
            \[
                Z=\frac{W}{\sigma\sqrt{\dfrac{\alpha^2}{n}+\dfrac{4\beta^2}{m}}}\sim\N(0,1).
            \]
            Por otro lado, en el denominador de $T$ podemos reescribir
            \[
                \frac{(n-1)S_X^2+\dfrac{(m-1)S_Y^2}{4}}{n+m-2}=\frac{\sigma^2U}{n+m-2},
            \]
            de modo que, sustituyendo sobre $T$ obtenemos
            \[
                T=\frac{W}{\sigma\sqrt{\dfrac{U}{n+m-2}}\;\sqrt{\dfrac{\alpha^2}{n}+\dfrac{4\beta^2}{m}}}
                =\frac{Z}{\sqrt{\dfrac{U}{n+m-2}}}.
            \]
            $Z$ depende solo de $(\overline{X},\overline{Y})$ y $U$ solo de $(S_X^2,S_Y^2)$, luego son independientes (por la extensión del Teorema de Fisher). Como $Z\sim\N(0,1)$ y $U\sim\chi^2_{n+m-2}$, por definición de la t de Student,
            \[
                \boxed{T\sim t_{n+m-2}}.
            \]

            \item Para $n=4$ y $m=5$ se tiene $n+m-2=7$, luego por el apartado anterior $T\sim t_7$. Consultando las tablas de la t de Student con $7$ grados de libertad, el valor que deja a su derecha una probabilidad de $0.025$ es
            \[
                t_{7;\,0.025}=2.365.
            \]
            Por tanto,
            \[
                \boxed{\P[T\geq 2.365]=0.025}.
            \]
        \end{enumerate}
    
        
    \end{ejercicio} 
\end{document}